This is a small CPU study with a single dataset of $8\times8$ images, a three-layer CNN and five seeds, so $p$-values are indicative and effects of about one standard deviation are not resolved. Baselines are reimplemented by us, not taken from the original code, and received no tuning: the supervised baseline uses a fixed number of steps without early stopping or validation, and the rotation baseline may be weak partly because of choices we did not explore. The test split is re-drawn per seed, so seeds are not independent samples of the population; the labeled subsets at $n=10$ (one image per class) make that budget very noisy. The sweeps and ablations are reported on the test split and were not used to choose the headline configuration, but a reader should treat any ``best'' value in them as post hoc. The unlabeled pool includes the labeled images (without labels), as is common in semi-supervised protocols. The PCA baseline has a fixed dimension (16) and standardised features, was not tuned, and is below raw pixels + LR at every budget, so pixels + LR is the stronger classical baseline here; the PCA gaps are thus likely inflated. We did not test other SSL methods (MoCo, BYOL), other probes, fine-tuning of the pretrained encoder, or larger images; a full study would add them with more seeds, a validation protocol and tuned baselines. Registry wall-clock times are inflated by machine contention. Finally, the gap to PCA may partly reflect the convolutional architecture rather than SSL, as the random-encoder control suggests.
